\chapter{Considerações parciais}
\label{chap_conclusao}

Esta dissertação, em sua versão de qualificação, investigou se, sob premissas técnicas de compatibilidade operacional, o biobunker pode ser considerado alternativa econômico-ambientalmente viável para apoiar a transição energética do transporte marítimo internacional operando a partir do Brasil, considerando as metas intermediárias de descarbonização da \gls{IMO} até 2030 e os condicionantes regulatórios associados a rotas internacionais. Para responder a essa questão, o trabalho desenvolveu modelo determinístico de simulação comparativa aplicado à rota Santos--Roterdã, articulando condições operacionais representativas, fatores de emissão em escopo \emph{well-to-wake}, custos de combustível, custo explícito de \emph{compliance} do \gls{EUETS} e critério de viabilidade baseado no custo marginal de abatimento de carbono.

À luz do critério de viabilidade adotado no modelo, a resposta à pergunta de pesquisa é negativa no cenário-base e positiva apenas de forma condicional em cenários regulatórios ampliados. No cenário-base e sob sinal regulatório compatível com a faixa recente do \gls{EUETS}, o biobunker não se mostra econômico-ambientalmente viável para nenhuma das rotas de produção avaliadas. Em todos os cenários simulados, o custo marginal de abatimento permanece acima do intervalo de referência adotado para o ambiente europeu de precificação de carbono. Isso significa que, embora a substituição parcial do combustível fóssil produza reduções de emissões, o ganho ambiental obtido não é suficiente, nas condições-base do modelo, para compensar o diferencial de custo da mistura frente ao patamar regulatório atualmente observado. Essa conclusão vale no escopo determinístico e paramétrico da modelagem, aplicado à rota Santos--Roterdã e sob os parâmetros médios adotados.

Ao mesmo tempo, os resultados mostram que essa resposta não deve ser interpretada de forma absoluta. A análise comparativa entre rotas de produção do biodiesel evidencia heterogeneidade relevante entre matérias-primas. Rotas associadas a \gls{UCO}, sebo bovino e \gls{TCO} apresentam desempenho superior ao da soja tanto em termos ambientais quanto em termos de preço de carbono de equilíbrio. Além disso, quando se ampliam conjuntamente os parâmetros regulatórios do modelo, passam a emergir cenários viáveis para \gls{UCO}, sebo bovino e \gls{TCO}, particularmente a partir de misturas em torno de \gls{B24}. Nesse sentido, os resultados parciais sustentam que o biobunker não é alternativa econômico-ambientalmente viável sob condições regulatórias correntes, mas passa a se mostrar econômico-ambientalmente viável nos cenários regulatórios ampliados simulados quando associado a matérias-primas de menor intensidade de carbono.

Esses achados permitem apresentar, já na qualificação, três contribuições principais do trabalho. A primeira consiste na proposição de um arcabouço analítico integrado capaz de avaliar, em um mesmo modelo, a articulação entre demanda útil constante, emissões \emph{WTW}, custo total por viagem e custo marginal de abatimento de carbono, evitando leituras isoladas apenas ambientais ou apenas econômicas. A segunda consiste na aplicação desse arcabouço a um caso representativo do transporte marítimo internacional operando a partir do Brasil, incorporando explicitamente a heterogeneidade entre rotas de produção do biodiesel e mostrando como diferenças de matéria-prima alteram o espaço de viabilidade. A terceira consiste na articulação entre modelagem quantitativa e interpretação regulatória, por meio da distinção entre custo explícito de \emph{compliance} e parâmetro de referência para avaliação da atratividade econômico-ambiental, o que permite ler a viabilidade do biobunker como resultado conjunto de custo, emissões e indução regulatória.

Como se trata ainda de etapa de qualificação, estas considerações têm caráter parcial e devem ser entendidas como base para aprofundamento posterior da dissertação. O estudo já permite delimitar com clareza o comportamento do modelo, as diferenças entre rotas e o papel do ambiente regulatório, mas também evidencia frentes concretas de continuidade metodológica e empírica. Entre elas, destacam-se o refinamento do domínio de mistura, o desacoplamento entre os sinais regulatórios do modelo, a expansão para novas rotas marítimas e o desenvolvimento de versões mais sensíveis à incerteza dos parâmetros econômicos, ambientais e operacionais.

Além dessas extensões de modelagem, a etapa seguinte do trabalho deverá priorizar a calibração empírica mais fina das rotas de biodiesel a partir de informações associadas ao contexto produtivo brasileiro. Isso envolve, por exemplo, o levantamento de dados mais precisos sobre preços, capacidade de oferta, distribuição regional de produção, matérias-primas predominantes e condições logísticas de abastecimento, com base em informações públicas de empresas brasileiras produtoras de biodiesel, associações setoriais, relatórios institucionais e bases regulatórias. Esse aprofundamento é relevante porque permitirá reduzir a dependência de parâmetros médios mais agregados e aproximar o modelo de condições mais realistas de formação de custo e disponibilidade das rotas analisadas.

No mesmo sentido, a continuidade da pesquisa deverá combinar ajustes no arquivo \texttt{simulacao.R} com nova rodada de análise dos resultados. Espera-se, nessa etapa, revisar a parametrização-base, atualizar sinais regulatórios e preços de referência, incorporar informações mais precisas sobre as rotas associadas a soja, sebo bovino, \gls{UCO} e \gls{TCO}, e avaliar se essas revisões alteram a fronteira de viabilidade identificada nesta versão. Paralelamente, a maturação do trabalho exigirá a consolidação desses achados em forma de artigo e o adensamento da redação final da dissertação, de modo a transformar os resultados parciais aqui apresentados em versão final consistente para defesa.

Com base nesse encaminhamento, a Tabela~\ref{tab:cronograma_qualificacao_defesa} apresenta o cronograma estratégico previsto do período posterior à qualificação até a defesa, tomando como referência o início das atividades em maio e o encerramento em dezembro.

\begin{table}[!ht]
\centering
\caption{Cronograma estratégico da qualificação à defesa}
\label{tab:cronograma_qualificacao_defesa}
\scriptsize{
\begin{tabular}{l c c c c c c c c}
\toprule
\multirow{2}{*}{Atividade} & \multicolumn{8}{c}{Meses} \\
& mai & jun & jul & ago & set & out & nov & dez \\
\midrule
Implementação & X & X & X & & & & & \\
Análise & & X & X & X & X & & & \\
Escrita de artigo & & & & X & X & X & & \\
Escrita da dissertação & & & X & X & X & X & X & \\
Defesa & & & & & & & & X \\
\bottomrule
\end{tabular}
}
\end{table}

Em termos estratégicos, a fase de implementação concentra os ajustes no modelo e na base empírica, incluindo a incorporação de parâmetros mais realistas para as rotas de biodiesel e eventuais modificações necessárias no \texttt{simulacao.R}. A fase de análise cobre a reexecução dos experimentos, a comparação entre resultados atualizados e os resultados parciais desta qualificação, bem como a interpretação das mudanças na fronteira de viabilidade. A escrita de artigo funciona como esforço de síntese e amadurecimento dos resultados centrais, enquanto a escrita da dissertação acompanha e consolida todo o processo até a preparação final para defesa.

Em síntese, esta etapa de qualificação já permite afirmar que o biobunker não se revela econômico-ambientalmente viável sob o ambiente regulatório corrente representado no cenário-base, mas tampouco deve ser tratado como alternativa irrelevante. Quando avaliado sob critérios explícitos de desempenho ambiental, custo e regulação, ele se apresenta como rota de transição condicional, cuja relevância depende da combinação concreta entre matéria-prima, mistura, contexto operacional e intensidade da indução regulatória. A etapa seguinte do trabalho deverá justamente aprofundar essa conclusão parcial, ampliando a base empírica do modelo e refinando os parâmetros que aproximam a simulação das condições efetivas do setor brasileiro de biodiesel e do transporte marítimo internacional.
